\documentclass[11pt,a4paper]{amsart}
\pdfinfoomitdate=1
\pdftrailerid{}
\pdfsuppressptexinfo=15
\usepackage[T1]{fontenc}
\usepackage{lmodern}
\usepackage[margin=27mm]{geometry}
\usepackage[expansion=false]{microtype}
\usepackage{amsmath,amssymb,mathtools}
\usepackage{array,booktabs}
\usepackage{xcolor}
\usepackage[numbers,sort&compress]{natbib}
\usepackage[colorlinks=true,linkcolor=blue!55!black,citecolor=blue!55!black,urlcolor=blue!55!black]{hyperref}
\usepackage[nameinlink,noabbrev]{cleveref}
\raggedbottom

\newtheorem{theorem}{Theorem}[section]
\newtheorem{proposition}[theorem]{Proposition}
\newtheorem{corollary}[theorem]{Corollary}
\newtheorem{lemma}[theorem]{Lemma}
\theoremstyle{definition}
\newtheorem{definition}[theorem]{Definition}
\theoremstyle{remark}
\newtheorem{remark}[theorem]{Remark}

\newcommand{\Tourn}{\mathcal T}
\newcommand{\Fix}{\operatorname{Fix}}
\newcommand{\score}{\operatorname{sc}}
\newcommand{\depth}{\operatorname{depth}}
\newcommand{\Energy}{\mathcal E}

\title[Score-upset reversal]{Synchronous Score-Upset Reversal in Tournaments:\
Energy, Recursive Blocks, and Fixed-Point Enumeration}
\author{Anonymous}
\date{}
\subjclass[2020]{05C20, 37P35, 05A15}
\keywords{tournament, score sequence, synchronous dynamics, ordinal sum, regular tournament, Artin--Mazur zeta function}

\begin{document}

\begin{abstract}
For a labelled tournament, simultaneously reorient every edge between
vertices of unequal outdegree from the higher-score endpoint to the
lower-score endpoint, while retaining edges whose endpoints have equal
score.  For this specified finite self-map, a strict squared-score Lyapunov
identity and a separate equal-score-block factorization prove termination.
The factorization identifies the pointwise hitting time with the height of a
well-founded score-refinement tree and gives the non-sharp universal bound
$\tau(T)\leq n-1$ for $n\geq1$.  It also identifies the fixed tournaments as
the unique ordered sums of regular tournaments.  Routine labelled
bookkeeping then gives, for regular-tournament counts $r_j$ and fixed counts
$f_n$,
\[
 f_0=1,\qquad f_n=\sum_{j=1}^n\binom nj r_jf_{n-j},\qquad
 \sum_{n\geq0}f_n\frac{x^n}{n!}
 =\frac{1}{1-\sum_{j\geq1}r_jx^j/j!}.
\]
All periodic points are fixed, so routine finite-map zeta bookkeeping gives
$(1-z)^{-f_n}$.  Score sequences, regular tournaments, ordinal sums,
regular-tournament enumeration, labelled exponential generating functions,
and zeta functions receive zero contribution credit.  After comparison with
static and iterative Copeland procedures and with arc/cycle-reversal work,
the residual scope is only the map-specific conjunction of the Lyapunov
identity, permanent block factorization, recursive depth and bound, and
fixed-set identification.  Ownership, priority, and external circulation
remain on hold.
\end{abstract}

\maketitle

\section{Introduction}

A tournament records the outcome of every pairwise contest.  Its most basic
local statistic is the score, or outdegree, of each vertex.  We study the
synchronous correction rule that points every unequal-score contest from
the currently higher-scoring vertex to the lower-scoring vertex and leaves
a score tie unchanged.  The rule looks like a one-pass ranking projection,
but it is not: changing contests also changes the next score classes.

Two simple structures control the resulting temporal feedback.  Reversing
a score upset strictly increases a quadratic score energy, even when many
reversals occur simultaneously.  Independently, the first update is an
ordinal sum over the old equal-score classes.  The score intervals of these
blocks cannot overlap, so every subsequent update runs separately inside
the unchanged induced subtournaments.  This second route gives a literal
recursive model of each orbit, not only a termination certificate.

The paper records four parts of one map-specific exact conjunction.
\begin{enumerate}
\item The squared-score energy has an exact positive increment on every
  nonfixed update; consequently, there are no nontrivial temporal cycles.
\item Equal-score blocks give an exact factorization of every positive
  iterate.  The pointwise depth equals the height of the associated
  refinement tree and is at most $n-1$.
\item Fixed tournaments are precisely unique ordered sums of regular
  tournaments.  A standard labelled-sequence corollary gives the recurrence
  and exponential generating function displayed in the abstract.
\item Every positive iterate has the same fixed set.  Standard finite-map
  zeta bookkeeping therefore reduces to one cycle factor.
\end{enumerate}

Tournament scores, their feasibility conditions, regular tournaments, and
score-based measures are classical topics
\cite{Landau1953,Moon1968,Monsuur2005}.  Static points rankings and Copeland
choice are likewise established \cite{Rubinstein1980,Henriet1985}, and
successive-choice procedures already iterate score-based selections
\cite{Bouyssou2004,LinaresBodanza2025}.  Score-preserving triangle reversals
and broader arc/cycle-reversal questions form another direct neighborhood
\cite{Ryser1964,Thomassen1988,GhoshEtAl2026}.  Regular-tournament enumeration,
labelled generating functions, and the periodic-point zeta construction are
also prior tools \cite{McKay1990,FlajoletSedgewick2009,ArtinMazur1965}.
Every one of these ingredients receives zero contribution credit here.  The
residual scope is only the conjunction of the exact statements for the
specific synchronous edge update.  A bounded owner search did not locate
that identical conjunction, but a search miss is neither novelty evidence
nor a priority certificate.  We make no claim about a sharp global depth or
a complete transient enumerator.  External dissemination remains
\textbf{HOLD}.

\section{Definition and the energy route}\label{sec:energy}

For $n\geq0$, put
\[
             [n]=\{i\in\mathbb Z:0\leq i<n\},
\]
so $[0]=\varnothing$ literally.  For a finite label set $V$, let
$\Tourn(V)$ be the set of tournaments on $V$, and put
$\Tourn_n=\Tourn([n])$.  Thus $\Tourn_0$ and $\Tourn_1$ are both singletons.
Write $u\to_Tv$ when $T\in\Tourn(V)$ directs the edge $\{u,v\}$ from $u$ to
$v$, and let
\[
 s_T(v)=d_T^+(v)
\]
be the score of $v$.

\begin{definition}[synchronous score-upset reversal]\label{def:update}
For $T\in\Tourn(V)$ and distinct $u,v\in V$, define
$\Phi_V(T)\in\Tourn(V)$ by
\begin{equation}\label{eq:update}
 u\to_{\Phi_V(T)}v
 \quad\Longleftrightarrow\quad
 s_T(u)>s_T(v)
 \quad\text{or}\quad
 \bigl(s_T(u)=s_T(v)\ \text{and}\ u\to_Tv\bigr).
\end{equation}
Every score on the right side is evaluated in the unchanged input $T$;
all edge decisions are then applied simultaneously.  We write
$\Phi_n=\Phi_{[n]}$ and suppress the subscript when the label set is clear.
For $C\subseteq V$, $T[C]$ denotes the induced tournament on $C$, and
$\Phi_C$ means this same definition on the label set $C$, using scores
internal to $T[C]$.

The hitting time is
\[
 \tau(T)=\min\{t\geq0:\Phi_V^{t+1}(T)=\Phi_V^t(T)\}.
\]
Before termination is proved, the minimum is understood as $\infty$ when
the displayed set is empty.
\end{definition}

Let
\[
 R(T)=\{(x,y)\in V^2:x\to_Ty\ \text{and}\ s_T(x)<s_T(y)\};
\]
these are exactly the arcs reversed by \eqref{eq:update}.  Put
\[
 \delta_v=s_{\Phi_V(T)}(v)-s_T(v),\qquad
 \Energy(T)=\sum_{v\in V}s_T(v)^2.
\]

\begin{theorem}[strict quadratic Lyapunov function]\label{thm:energy}
For every $T\in\Tourn(V)$,
\begin{equation}\label{eq:energy}
 \Energy(\Phi_V(T))-\Energy(T)
 =2\sum_{(x,y)\in R(T)}\bigl(s_T(y)-s_T(x)\bigr)
  +\sum_{v\in V}\delta_v^2.
\end{equation}
Moreover,
\begin{equation}\label{eq:local-fixed}
             \Phi_V(T)=T\quad\Longleftrightarrow\quad R(T)=\varnothing.
\end{equation}
Thus the energy difference is positive if and only if $\Phi_V(T)\ne T$.
Every orbit reaches a fixed point, and every periodic point is fixed.
\end{theorem}

\begin{proof}
The simultaneous score change has the incidence formula
\begin{equation}\label{eq:incidence}
 \delta_v=
 \bigl|\{x\in V:(x,v)\in R(T)\}\bigr|
 -\bigl|\{y\in V:(v,y)\in R(T)\}\bigr|.
\end{equation}
Expand the left side of \eqref{eq:energy} as
\[
 2\sum_{v\in V}s_T(v)\delta_v+\sum_{v\in V}\delta_v^2.
\]
By \eqref{eq:incidence},
\[
 \sum_{v\in V}s_T(v)\delta_v
 =\sum_{(x,y)\in R(T)}\bigl(s_T(y)-s_T(x)\bigr),
\]
so \eqref{eq:energy} follows without treating simultaneous reversals as
independent updates.  Each summand in its first term is a positive integer.
The definition reverses precisely the arcs in $R(T)$ and retains every
other arc, proving \eqref{eq:local-fixed} and strictness.

The phase space is finite.  Along a nonconstant orbit the integer
$\Energy$ strictly increases, so every orbit terminates and no temporal
cycle of length greater than one can occur.
\end{proof}

\section{Equal-score blocks and recursive depth}\label{sec:blocks}

For tournaments $S_1,\ldots,S_k$ on disjoint vertex sets, their
\emph{ordinal sum} $S_1\oplus\cdots\oplus S_k$ retains every internal edge
and directs every edge from $S_i$ to $S_j$ when $i<j$.

Fix $T\in\Tourn_n$.  Let $C_1,\ldots,C_k$ be its equal-score classes,
ordered so that their common scores decrease, and set
\[
 T_i=T[C_i],\qquad L_i=\sum_{j>i}|C_j|.
\]

\begin{theorem}[ordinal-sum and iterate factorization]\label{thm:factor}
The first update is
\begin{equation}\label{eq:first-factor}
       \Phi_n(T)=T_1\oplus T_2\oplus\cdots\oplus T_k.
\end{equation}
For $v\in C_i$,
\begin{equation}\label{eq:new-score}
       s_{\Phi_n(T)}(v)=L_i+s_{T_i}(v).
\end{equation}
In particular, the score intervals belonging to distinct old classes are
disjoint and retain their strict order.  For every $t\geq1$,
\begin{equation}\label{eq:iterate-factor}
 \boxed{\displaystyle
 \Phi_n^t(T)=
 \Phi_{C_1}^{\,t-1}(T_1)\oplus\cdots\oplus
 \Phi_{C_k}^{\,t-1}(T_k),}
\end{equation}
Here $T_i=T[C_i]$ throughout, and $\Phi_{C_i}$ is the restriction/update
operator defined before \cref{thm:energy}; in particular, its scores are
internal to $C_i$.
\end{theorem}

\begin{proof}
If $i<j$, every vertex of $C_i$ has larger old score than every vertex of
$C_j$.  The update therefore directs all edges from $C_i$ to $C_j$.  An
edge inside one class is tied and remains unchanged.  This proves
\eqref{eq:first-factor}.  A vertex in $C_i$ beats all $L_i$ vertices in
lower blocks, loses to all vertices in higher blocks, and retains its
internal wins.  This gives \eqref{eq:new-score}.

The scores of block $C_i$ lie in the integer interval
\[
             [L_i,L_i+|C_i|-1].
\]
The maximum of the next lower interval is $L_i-1$, so distinct block
intervals are disjoint.  Thus future updates never reverse an interblock
edge or merge two old classes.  Within $C_i$, every global score has the
same additive term $L_i$, so comparing global scores is the same as
comparing internal scores.  Applying this observation inductively proves
\eqref{eq:iterate-factor}.
\end{proof}

The factorization supplies a second, non-energy proof of convergence.  We
first justify the recursion it uses.  A tournament with one score class has
all scores tied, so \eqref{eq:update} retains every edge and the tournament
is fixed.  Consequently, a nonfixed tournament has at least two nonempty
score classes.  Each class is then a proper subset of the label set, so each
induced child has strictly smaller order.  The following recursion is
therefore well-founded: the score-refinement tree $\mathcal R(T)$ is a leaf,
with no children, when $T$ is fixed; otherwise its children are precisely
the trees $\mathcal R(T[C_i])$ of its score classes.  Let $h(T)$ be its
height, with a leaf of height zero.

\begin{corollary}[exact pointwise depth and universal bound]\label{cor:depth}
For every tournament,
\begin{equation}\label{eq:tree-depth}
                       \tau(T)=h(T).
\end{equation}
For $n\geq1$,
\begin{equation}\label{eq:depth-bound}
                       \tau(T)\leq n-1.
\end{equation}
The unique tournaments of orders zero and one have depth zero.
\end{corollary}

\begin{proof}
We use strong induction on the number of vertices.  The fixed case has
$\tau(T)=h(T)=0$.  If $T$ is nonfixed, every child has smaller order by the
well-foundedness argument, so the induction hypothesis gives finite
$\tau_i=h(T[C_i])$ for all children.
Since the blocks have fixed labels and all
interblock arcs are frozen, restriction to each $C_i$ shows that, for every
$q\geq1$,
\[
 \Phi_n^{q+1}(T)=\Phi_n^q(T)
 \quad\Longleftrightarrow\quad
 \Phi_{C_i}^{q}(T_i)=\Phi_{C_i}^{q-1}(T_i)
 \quad\text{for every }i.
\]
Equality at one time makes that factor fixed and hence persists at all later
times.  Thus all blocks are stable in the displayed equivalence if and only
if $q-1\geq\max_i\tau_i$.  Because the nonfixed parent is not stable at
$q=0$, this proves both directions of
\begin{equation}\label{eq:depth-recursion}
             \tau(T)=1+\max_i\tau(T[C_i]).
\end{equation}
This is exactly the recursive definition of $h(T)$.

For the bound, use induction on $n$.  A nonfixed tournament has at least two
score classes by the well-foundedness argument above.  Hence every $C_i$ has
size at most $n-1$.
Induction in \eqref{eq:depth-recursion} gives
\[
 \tau(T)\leq1+\max_i(|C_i|-1)\leq n-1.
\]
The cases $n=0,1$ follow directly from the definition.
\end{proof}

The proof of \cref{cor:depth} is structurally independent of
\cref{thm:energy}: it terminates the orbit by induction on block size and
also records the full pointwise recursion.  Conversely, the energy proof
excludes cycles without first identifying a score-class tree.

\section{Fixed points, enumeration, and zeta}\label{sec:fixed}

A tournament on $j$ vertices is \emph{regular} when every vertex has
outdegree $(j-1)/2$.  Nonempty regular tournaments therefore have odd
order.

\begin{theorem}[unique regular-block decomposition]\label{thm:fixed}
A tournament is fixed by score-upset reversal if and only if it is an
ordinal sum of nonempty regular tournaments.  This ordered decomposition is
unique: its blocks are the equal-score classes.  The empty tournament is
the empty ordinal sum.
\end{theorem}

\begin{proof}
Suppose $T$ is fixed and order its score classes as in
\cref{sec:blocks}.  By the local criterion \eqref{eq:local-fixed}, every
interclass edge points from the higher-score class to the lower-score class.
Thus $T=T[C_1]\oplus\cdots\oplus T[C_k]$.  Every vertex in $C_i$ has the
same $L_i$ external wins.  Since the vertices also have equal global
scores, their internal scores are equal, and $T[C_i]$ is regular.

Conversely, let $T=S_1\oplus\cdots\oplus S_k$ with regular block sizes
$m_1,\ldots,m_k$.  Every vertex in block $i$ has score
\[
 a_i=\sum_{j>i}m_j+\frac{m_i-1}{2}.
\]
For adjacent blocks,
$a_i-a_{i+1}=(m_i+m_{i+1})/2>0$.  Hence each block is one score class,
all interclass edges already point downward in score, and tied internal
edges remain unchanged.  The decomposition is unique because score classes
are intrinsic to $T$.
\end{proof}

Let $r_j$ be the number of regular tournaments on a fixed labelled
$j$-set.  We use $r_0=0$ because a block is required to be nonempty, while
$r_j=0$ for positive even $j$.  Let
\[
                 f_n=|\Fix(\Phi_n)|,
\]
where $f_0=1$ counts the empty ordered sum.

\begin{corollary}[zero-credit labelled bookkeeping]
\label{cor:enumeration}
For $n\geq1$,
\begin{equation}\label{eq:fixed-rec}
 \boxed{\displaystyle
 f_n=\sum_{j=1}^n\binom nj r_jf_{n-j}.}
\end{equation}
With
\[
 R(x)=\sum_{j\geq1}r_j\frac{x^j}{j!},\qquad
 F(x)=\sum_{n\geq0}f_n\frac{x^n}{n!},
\]
one has the formal identity
\begin{equation}\label{eq:egf}
                       \boxed{F(x)=\frac{1}{1-R(x)}.}
\end{equation}
The first values are
\[
 (f_0,f_1,\ldots,f_6)=(1,1,2,8,40,264,2048).
\]
\end{corollary}

\begin{proof}
For a nonempty fixed tournament, choose its unique highest-score regular
block.  If its size is $j$, there are $\binom nj$ choices of labels,
$r_j$ internal orientations, and $f_{n-j}$ choices for the ordered sum
below it.  This proves \eqref{eq:fixed-rec}.  Multiplying by $x^n/n!$ and
summing gives $F(x)=1+R(x)F(x)$, which is equivalent to
\eqref{eq:egf}; this is the standard labelled-sequence construction
\cite{FlajoletSedgewick2009}.  The displayed values follow from
$(r_1,r_3,r_5)=(1,2,24)$, part of the prior regular-tournament enumeration
literature \cite{McKay1990}.
\end{proof}

For a self-map $G$ of a finite set, write
\[
 \zeta_G(z)=\exp\left(\sum_{m\geq1}
                \frac{|\Fix(G^m)|}{m}z^m\right).
\]

\begin{corollary}[zero-credit periodic census and zeta]\label{cor:zeta}
For every $n\geq0$ and $m\geq1$,
\[
             \Fix(\Phi_n^m)=\Fix(\Phi_n),
\]
and therefore
\begin{equation}\label{eq:zeta}
                 \boxed{\zeta_{\Phi_n}(z)=(1-z)^{-f_n}.}
\end{equation}
In particular, the boundary systems $n=0,1$ both have zeta function
$(1-z)^{-1}$.
\end{corollary}

\begin{proof}
Every fixed point of $\Phi_n$ is fixed by every positive iterate.  The
reverse inclusion follows from \cref{thm:energy}, which excludes temporal
cycles of length greater than one.  Substitution in the defining series and
$\sum_{m\geq1}z^m/m=-\log(1-z)$ give \eqref{eq:zeta}; the definition itself
is classical \cite{ArtinMazur1965}.
\end{proof}

\section{Exact controls, owner subtraction, and scope}\label{sec:controls}

The accompanying standard-library verifier exhausts every labelled
tournament for $0\leq n\leq6$, a total of $33{,}868$ states.  It compares
bit-coded and literal set-of-arcs updates, checks the exact energy identity,
the score-interval formula, every factorized iterate through time $n$, the
tree-depth equality, the regular-block criterion, the recurrence, and fixed
sets of the first eight positive iterates.  Its $1{,}677{,}508$ assertions
all pass.  Selected complete lanes are
\begin{center}
\begin{tabular}{@{}rrrrl@{}}
\toprule
$n$ & $|\Tourn_n|$ & fixed & maximum observed depth & depth histogram\\
\midrule
0 & 1      & 1     & 0 & $\{0:1\}$\\
1 & 1      & 1     & 0 & $\{0:1\}$\\
2 & 2      & 2     & 0 & $\{0:2\}$\\
3 & 8      & 8     & 0 & $\{0:8\}$\\
4 & 64     & 40    & 1 & $\{0:40,1:24\}$\\
5 & 1,024  & 264   & 1 & $\{0:264,1:760\}$\\
6 & 32,768 & 2,048 & 2 & $\{0:2048,1:26400,2:4320\}$\\
\bottomrule
\end{tabular}
\end{center}

The same exhaustion actively searches for idempotence failure.  Order the
edges lexicographically as
\[
 (0,1),(0,2),\ldots,(0,5),(1,2),\ldots,(4,5),
\]
and let bit one mean that the smaller endpoint wins.  The six-vertex mask
$148$ has the exact orbit
\[
 148\;\xmapsto{\Phi_6}\;4\;\xmapsto{\Phi_6}\;0
 \;\xmapsto{\Phi_6}\;0
\]
and score vectors
\[
 (2,2,2,2,3,4),\qquad(1,1,2,2,4,5),\qquad(0,1,2,3,4,5).
\]
Thus the system is not a projection.  More precisely, the program scans
$n=0,1,\ldots,6$ increasingly and, at each order, scans the numerical masks
in increasing order.  Mask $148$ is the least mask with
$\Phi_n^2(T)\ne\Phi_n(T)$ in this specified finite scan; no lower scanned
order has such a state.  This order-dependent statement is the entire
minimality claim.  It neither identifies a sharp global depth nor replaces
the infinite proofs.

\paragraph{Mechanics of the closest procedure families.}
Every row in the comparison below receives zero contribution credit.
\begin{center}
\scriptsize
\begin{tabular}{@{}
  >{\raggedright\arraybackslash}p{.22\linewidth}
  >{\raggedright\arraybackslash}p{.33\linewidth}
  >{\raggedright\arraybackslash}p{.37\linewidth}@{}}
\toprule
work & state/output and deletion & arc action, score timing, and ties\\
\midrule
Rubinstein; Henriet
 & static ranking or choice from a tournament; the input is not advanced
 & no arc changes and no temporal tie-retention rule; scores are evaluated
   for a static output \cite{Rubinstein1980,Henriet1985}\\
Bouyssou
 & ranking by successive choice; selected alternatives are removed
 & no arc reorientation; stages are sequential and the choice rule is
   recomputed on a shrinking set \cite{Bouyssou2004}\\
Linares Lejarraga--Bodanza
 & an iterative collective-choice procedure, not an orientation-valued
   self-map
 & no synchronous all-edge correction from one old score vector; its
   ``Iterative Copeland from below'' vocabulary is not our update
   \cite{LinaresBodanza2025}\\
Ryser; Thomassen; Ghosh et al.
 & tournament orientations; no vertex deletion
 & chosen triangle/cycle or arc reversals, sequential or controlled;
   Ryser's and the ESA score-sequence setting is score-preserving, rather
   than the present simultaneous score-changing rule
   \cite{Ryser1964,Thomassen1988,GhoshEtAl2026}\\
present map
 & the complete orientation is the state; no vertex deletion
 & all unequal-score pairs are oriented simultaneously from one unchanged
   old score vector; tied edges are retained; scores are recomputed only for
   the next time step\\
\bottomrule
\end{tabular}
\end{center}

\paragraph{Zero-credit owner ledger.}
Landau's score-sequence theorem receives zero credit
\cite{Landau1953}.  Moon's tournament terminology, regular-tournament and
ordinal-sum background, including the route to Ryser's triangle-reversal
theorem, receives zero credit \cite{Moon1968,Ryser1964}.  Static score
ranking and Copeland choice receive zero credit
\cite{Rubinstein1980,Henriet1985}; successive and current iterative
Copeland-family choice procedures receive zero credit
\cite{Bouyssou2004,LinaresBodanza2025}.  Score-upset and inconsistency
language receives zero credit \cite{Monsuur2005}.  The arc/cycle-reversal
lineage, including the contemporary ESA 2026 score-sequence neighbor,
receives zero credit \cite{Thomassen1988,GhoshEtAl2026}.  Regular-tournament
counts receive zero credit \cite{McKay1990}; the generic labelled-sequence
recurrence and exponential generating function receive zero credit
\cite{FlajoletSedgewick2009}; and the zeta definition and its one-factor
specialization receive zero credit \cite{ArtinMazur1965}.

After these subtractions, the only residual scope under consideration is
the exact finite-map conjunction for \eqref{eq:update}: the simultaneous
incidence/Lyapunov identity, permanent score-class factorization, recursive
pointwise depth formula with its universal upper bound, and the map-specific
identification of the fixed set.  The recurrence, exponential generating
function, regular counts, and zeta line are low-credit corollaries, not
independent contributions.  The bounded audit covered tournament score
correction, synchronous/parallel score updates, Copeland iteration and
ranking by choosing, arc/cycle reversal, and regular-tournament enumeration.
Its failure to locate an identical update is not evidence of novelty or
priority.  Discovery of a temporal owner would supersede the residual scope;
owner clearance remains \textbf{HOLD}.

\paragraph{Internal P106 firewall.}
The P106 MIS-polarity system evolves subsets of the vertex set of a fixed
undirected graph by an antitone neighborhood operator.  Its proof engine is
a Galois polarity and its recurrent objects are independent dominating
sets.  Here the state is the entire orientation of a complete graph, every
update may change edges, regular cyclic tournaments can be fixed, and the
proof engines are score energy and ordinal-sum refinement.  Shared words
such as ``synchronous,'' ``fixed point,'' and ``zeta'' receive no
contribution credit; neither phase space nor update is a parameter variant
of P106.

\paragraph{Limits.}
We do not determine the sharp maximum of $\tau$ as a function of $n$, and
we do not enumerate all transient tree heights.  The bound $n-1$ is the
proved universal statement.  The bounded owner search is not an external
literature certificate.  The anonymous working manuscript remains under
external \textbf{HOLD}.

\section{Conclusion}

Synchronous correction of score upsets has more structure than a one-pass
ranking operation.  A strict quadratic energy rules out recurrence, while
disjoint score intervals turn every orbit into a recursive refinement of
unchanged induced subtournaments.  That recursion gives the exact
pointwise hitting time and the universal $n-1$ bound.  At equilibrium, the
same score classes become the unique regular blocks, which yields the fixed
recurrence, its exponential generating function, and the one-factor zeta
formula as zero-credit bookkeeping consequences.  The six-vertex state from
the specified finite scan shows why a recursive step can be necessary.

\bibliographystyle{plainnat}
\bibliography{references}

\end{document}
